\bumppoints{34}

\question
Based on the hypothetical food web below, one letter represents an organism that could be both a secondary and tertiary consumer? The arrows point from lower to higher trophic levels. More than one answer is possible. Circle only one answer.

\begin{multicols}{3} 
	\begin{choices}
\correctchoice A %dontshuffle 
\choice B %dontshuffle 
\choice C %dontshuffle 
\choice D %dontshuffle 
\correctchoice E %dontshuffle 
\end{choices}
	\columnbreak
	\includegraphics{exam4_foodweb}
	\columnbreak
\end{multicols}
\vspace{-1\baselineskip}

\question
The dead zone in the Gulf of Mexico and other oceans around the globe occurs when

\begin{choices}
\choice eutrophication causes an absence of nutrients in the zone.
\choice global climate change causes the water to become too warm to sustain life.
\choice commercial fishers over-harvest shrimp, crab, fishes and other sea animals.
\correctchoice an excess of nutrients causes an algal bloom, followed by death of the algae. Decomposition of the algae uses up the oxygen in the water, so other organisms can no longer live there.
\choice an excess of nutrients causes an algal bloom. The algae consume all of the oxygen in the water for photosynthesis, so other organisms can no longer live there.
\end{choices}

\question
The Ozark region of Missouri lies within the temperate deciduous forest ecosystem.  If, within the next 25 years, global climate change raises the mean annual temperature from 14°C to 18°C and lowers the mean annual precipitation from 160 cm to 130 cm, which of the following is most likely?

\begin{multicols}{2}
	\begin{choices}
		\choice The ecosystem will become more like temperate grassland. 
		\correctchoice The ecosystem will become more like sclerophyllous or thorn woodland. 
		\choice The ecosystem will become more like temperate rain forest. 
		\choice The ecosystem will remain temperature deciduous ecosystem. 
		\choice The ecosystem will become more like boreal forest. 
	\end{choices}
	\columnbreak
	\includegraphics[width=0.45\textwidth]{exam4_climograph2}
\end{multicols}

\question
Considering your knowledge of the dead zone and the causes of microscopic algae blooms, the Gulf of Mexico ecosystem is \emph{most likely} to be limited by which of the following nutrients?

\begin{choices}
\choice oxygen
\choice carbon
\choice hydrogen
\correctchoice nitrogen
\choice \ce{CO2}
\end{choices}

\question
Plants produce water vapor during photosynthesis, which goes from the plant to the atmosphere. The water vapor in the atmosphere eventually falls back to the earth as rain. This process would be part of

\begin{choices}
\correctchoice the water cycle.
\choice the nutrient cycle.
\choice energy flow.
\choice the energy cycle.
\choice plant metabolism.
\choice the entire ecosystem.
\end{choices}

\question
Which of these species is most likely a keystone species in Alaska?

\begin{choices}
\correctchoice Sea otters. 
\choice Kelp.  
\choice Humans.  
\choice Killer whales. 
\choice Sea urchins.  
\end{choices}

\question
A tertiary consumer has high abundance. In a trophic cascade, which of the following is probably \emph{not} true.
\begin{choices}
\choice The secondary consumer will have low abundance. 
\correctchoice The primary producer will have high abundance. 
\choice The decomposer level will not be affected. 
\choice The primary consumer will have high abundance. 
\end{choices}

\question
If humans stopped fishing along the coast of Alaska, what will probably happen to the kelp forest and why?

\begin{choices}
\choice The kelp abundance will decrease because more fish will be present to eat the kelp.
\choice The kelp abundance will decrease. Sea otters eat urchins and urchins eat kelp. Because there were fewer otters to eat urchins, the number of urchins increased. The urchins therefore ate more kelp so the kelp abundance decreased.
\correctchoice The kelp abundance will increase because the sea urchin abundance will probably decrease.
\choice The kelp abundance will increase because of the symbiotic relationship between kelp and fishes. Fishes are necessary for kelp to reproduce.
\choice The kelp abundance will remain about the same. The kelp are not part of this community, so the kelp would not be affected.
\end{choices}

\question
In ecosystems, why is ``cycling'' used to describe the transfer of chemical elements and nutrients, while ``flow'' is used for the transfer of energy?

\begin{choices}
	\choice Chemical elements are cycled into ecosystems from other ecosystems, but energy constantly flows within the ecosystem.
	\choice Both chemical elements and energy flow within an ecosystem but can cycle to other ecosystems.
	\choice Energy flows from one ecosystem to other ecosystems, but chemical elements constantly cycle within an ecosystem.
	\choice Both chemical elements and energy are recycled but can flow to other ecosystems.
	\correctchoice Chemical elements are recycled within ecosystems, but energy flows through and out of ecosystems.
\end{choices}

\question
The Mediterranean coast and coastal California have similar types of ecosystems. This is because

\begin{choices}
\choice they are at similar latitudes.
\choice they are both wine-growing regions.
\correctchoice they have similar climates.
\choice they have similar types of organisms.
\choice they have similar types of communities.
\end{choices}

\question
Which block in the figure below shows the greatest species diversity? (Different shades, including white, represent different species. The small squares indicate relative abundance of individuals present.)

\begin{multicols}{3}
	\begin{choices}
		\choice \includegraphics{diversity_mcC} 
		\choice \includegraphics{diversity_mcB} 
		\choice \includegraphics{diversity_mcA} 
		\choice \includegraphics{diversity_mcD} 
		\correctchoice \includegraphics{diversity_mcE} 
	\end{choices}
\end{multicols}

\question
In 2011, the Mississippi River experienced a record flood event.  Which of the following statements is the best prediction of how this record flood event will affected the Gulf of Mexico dead zone that summer?

\begin{choices}
	\correctchoice The dead zone will be larger because the volume of water will carry even more nutrients into the Gulf of Mexico.
	\choice The dead zone will be larger because the volume of water will cause more fishes to avoid the freshwater input.
	\choice The dead zone will be smaller because the volume of water will dilute the nutrients, reducing eutrophication.
	\choice The dead zone will be smaller because the volume of water will bring more oxygen to the Gulf of Mexico, which will offset the hypoxia.
	\choice The dead zone will be similar to previous years because most of the flood water will soak into the ground and not enter the Gulf of Mexico.
\end{choices}

\question
The energy stored by producers can be passed through an ecosystem along a(n) \rule{1in}{0.4pt}, a series of interactions in which organisms transfer energy by eating and being eaten.

\begin{choices}
	\choice predator-prey chain  
	\choice energy cycle 
	\correctchoice food web 
	\choice trophic cascade 
	\choice ecosystem 
	\choice nutrient cycle 
\end{choices}

\question
A fish ate shrimp. The shrimp ate algae (which are photosynthetic).  The algae is a \rule{1in}{0.4pt}. The fish is a \rule{1in}{0.4pt}.

\begin{choices}
	\choice primary producer; primary consumer. 
	\choice primary consumer; primary producer. 
	\choice secondary consumer; primary consumer. 
	\correctchoice primary producer; secondary consumer. 
	\choice secondary consumer; primary producer. 
	\choice primary consumer; secondary consumer. 
\end{choices}

\question
The land on the south side of the Himalayan mountain range (below) receives relatively high levels of rainfall and has fertile farmlands.  The land north of the mountain range is hot and dry and can't be farmed.  Which of the following is true.

\begin{multicols}{2}
\begin{choices}
\choice The south side of the Himalayan mountain range is in a rain shadow.
\choice The warm, moist air is too dense to rise over the height of Mount Everest and the other mountains in this range so the moisture remains on the south side of the range.
\correctchoice Warm, moist air from the Indian Ocean blows north, rises and cools, releasing precipitation as it moves up the south side of the range, and this cool, now dry air mass heats up as it descends on the north side of the range.
\choice The desert in southern China warms the north side of the mountain range, causing it to dry out.
\choice The cool, moist Indian Ocean air heats up as it rises, releasing its precipitation as it passes the tops of the mountains, and this warm, now dry air cools as it descends on the leeward side of the range.
\end{choices}
\columnbreak
	\includegraphics[width=0.45\textwidth]{exam4_himalayas}
\end{multicols}

\newpage

\question
Consider the food chain grass $\to$ grasshopper $\to$ mouse $\to$ snake $\to$ hawk. How much of the chemical energy fixed by photosynthesis of the grass (100\%) is available to the hawk?

\begin{choices}
	\correctchoice 0.01\%
	\choice 0.1\% %dontshuffle 
	\choice 1\% %dontshuffle 
	\choice 10\% %dontshuffle 
	\choice 60\% %dontshuffle 
\end{choices}

\question
On average, the amount of secondary consumer biomass is about \underline{\hspace{2cm}} less than the amount of primary producer biomass. 
\begin{choices}
	\choice 10$\times$ %dontshuffle 
	\choice 20$\times$ %dontshuffle 
	\choice 50$\times$ %dontshuffle 
	\choice 80$\times$ %dontshuffle 
	\correctchoice 100$\times$%dontshuffle 
\end{choices}


